%! AP Statistics — Unit 1, Lesson 1.10 — Guided Notes
\documentclass[10pt]{article}
\usepackage{apstats-article}
\usepackage{apstats-boxes}

%=============================================================================
\begin{document}

\pageheader{Unit 1, Lesson 1.10}{Guided Notes}
\namedateperiod

\vspace{0.10in}

% ── OBJECTIVE ────────────────────────────────────────────────────────────────
\begin{objectivebox}
\small By the end of this lesson, I will be able to\dots\\[4pt]
\writeline\\[1pt]
\writeline\\[1pt]
\writeline
\end{objectivebox}

\vspace{0.10in}

% ── VOCABULARY ───────────────────────────────────────────────────────────────
\begin{vocabbox}
\begin{multicols}{2}
\termblanklong{Normal distribution}
\termblanklong{68--95--99.7 Rule (Empirical Rule)}
\termblanklong{$z$-score (standardized value)}
\columnbreak
\termblanklong{Standard Normal distribution $N(0,1)$}
\termblanklong{Normal table (Table A)}
\termblanklong{Normal probability plot (NPP)}
\end{multicols}
\end{vocabbox}

\vspace{0.10in}

% ── HOOK ─────────────────────────────────────────────────────────────────────
\begin{hookbox}
\small\textbf{Why does the Normal model matter?}

\vspace{5pt}
Adult male heights follow approximately $N(70, 3)$ inches. Without any data, we can
answer: \textit{``What percent of men are between 67 and 73 inches tall?''} --- because
we know the distribution's shape and parameters. This is the power of a mathematical model.

\vspace{4pt}
Today's question: if we know a distribution is approximately Normal,
\textbf{what proportion} falls in any region we choose?
\end{hookbox}

\vspace{0.08in}

% ── SECTION 1: SHAPE AND PARAMETERS ──────────────────────────────────────────
\begin{notesbox}{Shape and Parameters of the Normal Distribution}
\small
\begin{multicols}{2}

\textbf{\textcolor{navy}{Properties of a Normal distribution:}}
\begin{itemize}[itemsep=3pt]
  \item Shape: \blank{3.5cm} (bell-shaped)
  \item Symmetric about the \blank{2cm}
  \item Mean $=$ Median $=$ Mode
  \item Described by \blank{1.5cm} parameters:
        $\mu$ (\blank{1.8cm}) and $\sigma$ (\blank{2.5cm})
  \item Notation: \blank{3cm}
  \item Tails extend to $\pm\infty$ but curve
        approaches \blank{2cm} quickly
\end{itemize}

\vspace{6pt}
\textbf{\textcolor{navy}{Changing $\mu$ shifts the curve:}}\\
Moving $\mu$ left or right \blank{3cm} the shape.

\vspace{4pt}
\textbf{\textcolor{navy}{Changing $\sigma$ rescales the curve:}}\\
Larger $\sigma \Rightarrow$ \blank{3cm} and \blank{2cm}.\\
Smaller $\sigma \Rightarrow$ \blank{3cm} and \blank{2cm}.

\columnbreak

\textbf{\textcolor{navy}{Pre-drawn Normal curve — $N(70, 3)$:}}

\vspace{4pt}
\begin{center}
\begin{tikzpicture}[scale=0.80]
  \draw[thick, navy]
    plot[domain=-3.5:3.5, samples=80] (\x, {2.5*exp(-\x*\x/2)});
  \draw[thick] (-3.8, 0) -- (3.8, 0);
  \foreach \z/\h in {-3/61, -2/64, -1/67, 0/70, 1/73, 2/76, 3/79} {
    \draw (\z, -0.12) -- (\z, 0.12);
    \node[below, font=\tiny] at (\z, -0.15) {\h};
  }
  \node[below, font=\tiny\itshape] at (0, -0.52) {Height (in.)};
  \node[above, font=\tiny\bfseries\color{navy}] at (0, 2.52) {$\mu = 70$};
  % SD labels
  \node[font=\tiny, color=charcoal] at (-1.5, 1.0) {$\sigma\!=\!3$};
  \draw[->, thin, charcoal] (-1.55, 0.92) -- (-1.0, 0.30);
\end{tikzpicture}
\end{center}

\vspace{4pt}
\small The curve is \blank{3cm} at $\mu = 70$
and falls symmetrically on both sides.

\end{multicols}
\end{notesbox}

\vspace{0.08in}

% ── SECTION 2: EMPIRICAL RULE ─────────────────────────────────────────────────
\begin{notesbox}{The 68--95--99.7 (Empirical) Rule}
\small
\begin{multicols}{2}

\textbf{\textcolor{navy}{The Rule:}} In any Normal distribution $N(\mu, \sigma)$:
\begin{itemize}[itemsep=4pt]
  \item $\approx \blank{2cm}\%$ of values fall
        within $\mathbf{1}\,\sigma$ of the mean
        \hfill $(\mu - \sigma$ to $\mu + \sigma)$
  \item $\approx \blank{2cm}\%$ of values fall
        within $\mathbf{2}\,\sigma$ of the mean
        \hfill $(\mu - 2\sigma$ to $\mu + 2\sigma)$
  \item $\approx \blank{2cm}\%$ of values fall
        within $\mathbf{3}\,\sigma$ of the mean
        \hfill $(\mu - 3\sigma$ to $\mu + 3\sigma)$
\end{itemize}

\vspace{6pt}
\textbf{\textcolor{navy}{Heights $N(70,3)$ — apply the rule:}}
\begin{itemize}[itemsep=3pt]
  \item 68\%: between \blank{1.5cm} and \blank{1.5cm} in.
  \item 95\%: between \blank{1.5cm} and \blank{1.5cm} in.
  \item 99.7\%: between \blank{1.5cm} and \blank{1.5cm} in.
\end{itemize}

\vspace{6pt}
\textbf{\textcolor{navy}{Finding a tail:}}\\
$P(\text{height} > 76) = \dfrac{100\% - 95\%}{2} = \blank{2cm}$

\columnbreak

\textbf{\textcolor{navy}{Empirical Rule diagram:}}

\vspace{4pt}
\begin{center}
\begin{tikzpicture}[scale=0.85]
  % Draw nested shaded regions (outermost first)
  \fill[navy!10]
    plot[domain=-3:3, samples=60] (\x, {2.5*exp(-\x*\x/2)})
    -- (3,0) -- (-3,0) -- cycle;
  \fill[navy!20]
    plot[domain=-2:2, samples=60] (\x, {2.5*exp(-\x*\x/2)})
    -- (2,0) -- (-2,0) -- cycle;
  \fill[sky]
    plot[domain=-1:1, samples=40] (\x, {2.5*exp(-\x*\x/2)})
    -- (1,0) -- (-1,0) -- cycle;
  \draw[thick, navy]
    plot[domain=-3.5:3.5, samples=80] (\x, {2.5*exp(-\x*\x/2)});
  % Vertical reference lines
  \foreach \z in {-3,-2,-1,1,2,3} {
    \draw[dashed, charcoal, thin] (\z, 0) -- (\z, {2.5*exp(-\z*\z/2)});
  }
  % Axis
  \draw[thick] (-3.8, 0) -- (3.8, 0);
  \foreach \z in {-3,-2,-1,0,1,2,3} {
    \draw (\z, -0.12) -- (\z, 0.12);
    \node[below, font=\tiny] at (\z, -0.15) {$\z$};
  }
  \node[below, font=\tiny\itshape] at (0, -0.52) {$z$-score};
  % Percent labels
  \node[font=\tiny\bfseries] at (0, 1.0) {68\%};
  \node[font=\tiny\bfseries] at (1.55, 0.45) {95\%};
  \node[font=\tiny\bfseries] at (2.55, 0.15) {99.7\%};
\end{tikzpicture}
\end{center}

\end{multicols}
\end{notesbox}

\vspace{0.08in}

% ── SECTION 3: z-SCORES AND NORMAL TABLE ──────────────────────────────────────
\begin{notesbox}{$z$-Scores and the Normal Table}
\small
\begin{multicols}{2}

\textbf{\textcolor{navy}{$z$-score formula:}}
\[
  z = \dfrac{x - \mu}{\sigma}
\]
Interpretation: $z$ tells us how many \blank{3cm}
above or below the \blank{1.5cm} the value $x$ falls.

\vspace{5pt}
\textbf{\textcolor{navy}{Standardizing:}} Converting $x$ to $z$ transforms
any $N(\mu, \sigma)$ into the \blank{3cm}
Normal distribution: $N(\blank{1cm},\blank{1cm})$.

\vspace{5pt}
\textbf{\textcolor{navy}{Example (heights $N(70,3)$):}}\\
$x = 76$: $z = \dfrac{76-70}{3} = \blank{1.5cm}$\\[4pt]
$x = 64$: $z = \dfrac{64-70}{3} = \blank{1.5cm}$\\[4pt]
$x = 70$: $z = \dfrac{70-70}{3} = \blank{1.5cm}$

\columnbreak

\textbf{\textcolor{navy}{Using Table A (Normal table):}}\\
Table A gives $P(Z < z)$ = area to the \blank{1.5cm}.

\vspace{4pt}
\renewcommand{\arraystretch}{1.5}
\begin{tabular}{@{}l l@{}}
$P(Z < z)$        & look up $z$ directly \\
$P(Z > z)$        & $= 1 - P(Z < z)$ \\
$P(z_1 < Z < z_2)$ & $= P(Z < z_2) - P(Z < z_1)$ \\
\end{tabular}

\vspace{6pt}
\textbf{\textcolor{navy}{Calculator:}}\\
\texttt{normalcdf(low, high, $\mu$, $\sigma$)}\\
\texttt{invNorm(area, $\mu$, $\sigma$)}

\vspace{6pt}
\textbf{\textcolor{navy}{Inverse Normal (percentile $\to$ value):}}\\
Find $z^*$ such that $P(Z < z^*) = p$, then\\
$x = \mu + z^* \cdot \sigma$

\end{multicols}
\end{notesbox}

\vspace{0.08in}

% ── SECTION 4: ASSESSING NORMALITY ───────────────────────────────────────────
\begin{notesbox}{Assessing Normality}
\small
\begin{multicols}{2}

\textbf{\textcolor{navy}{Two graphical methods:}}
\begin{enumerate}[itemsep=3pt, label=\alph*.]
  \item \textbf{Histogram} (or dotplot): look for a
        \blank{3cm}, \blank{2cm} shape with
        no strong \blank{2cm}.
  \item \textbf{Normal probability plot (NPP)}:
        if the points fall \blank{3cm} along a
        diagonal line, the distribution is approximately
        Normal. \blank{2cm} from a line suggest skewness.
\end{enumerate}

\vspace{5pt}
\textbf{\textcolor{navy}{When can we use the Normal model?}}\\
Use it when: the histogram is \blank{3cm} and
\blank{3cm}, or the NPP is \blank{2.5cm}. Never
use it blindly without \blank{2.5cm} first.

\columnbreak

\textbf{\textcolor{navy}{Normal probability plot examples:}}

\vspace{4pt}
\begin{center}
\begin{tikzpicture}[scale=0.7]
  % Approximately Normal NPP (roughly linear)
  \draw[thick] (0,0) -- (3.5, 3.5);
  \foreach \x/\y in {0.3/0.25, 0.8/0.85, 1.3/1.25, 1.8/1.82, 2.3/2.28, 2.8/2.75, 3.2/3.18} {
    \filldraw[navy] (\x, \y) circle (0.07);
  }
  \node[below, font=\tiny\itshape] at (1.75, -0.15) {Approx.\ Normal};
  \draw[thin, charcoal] (0,0) rectangle (3.5, 3.5);
\end{tikzpicture}
\hspace{0.5cm}
\begin{tikzpicture}[scale=0.7]
  % Non-Normal NPP (curved)
  \draw[thick] (0,0) -- (3.5, 3.5);
  \foreach \x/\y in {0.3/0.05, 0.8/0.30, 1.3/0.85, 1.8/1.80, 2.3/2.75, 2.8/3.20, 3.2/3.45} {
    \filldraw[navy] (\x, \y) circle (0.07);
  }
  \node[below, font=\tiny\itshape] at (1.75, -0.15) {Not Normal (skewed)};
  \draw[thin, charcoal] (0,0) rectangle (3.5, 3.5);
\end{tikzpicture}
\end{center}

\vspace{4pt}
\small
Left plot: points near the line $\Rightarrow$ \blank{2.5cm} Normal.\\
Right plot: points curve away $\Rightarrow$ \blank{2.5cm}.

\end{multicols}
\end{notesbox}

\vspace{0.08in}

% ── QUICK CHECK ──────────────────────────────────────────────────────────────
\begin{tcolorbox}[colback=goldbg, colframe=goldacc, boxrule=0.8pt, arc=2mm,
    left=4mm, right=4mm, top=2mm, bottom=2mm]
\small\textbf{Quick Check:} Newborn weights are approximately $N(7.5, 1.1)$ pounds.
(a) What proportion of newborns weigh between 6.4 and 8.6 lb? Use the Empirical Rule.
\hfill $P = \blank{2cm}$\\[4pt]
(b) What is the $z$-score for a newborn weighing 9.7 lb?
\hfill $z = \blank{2cm}$\\[4pt]
(c) Is 9.7 lb unusually high? \blank{2cm} \quad Justify: \blank{6cm}
\end{tcolorbox}

\end{document}
